% Part:sets-functions-relations
% Chapter: size-of-sets
% Section: equinumerous-sets

\documentclass[../../../include/open-logic-section]{subfiles}

\begin{document}

\olfileid[tr]{sfr}{siz}{equ}

\olsection{Eş Güçlülük}

Sonlu kümeler için iyi işleyen sezgisel bir küme ``büyüklüğü'' kavramımız
var. Peki sonsuz kümeler ne olacak? \emph{Herhangi} bir büyüklükteki iki
kümenin büyüklüğünü biçimsel olarak karşılaştırmak istiyorsak işe kümelerin
ne zaman aynı büyüklükte olduğunu tanımlayarak başlamak iyi bir fikirdir.
Frege bunu şöyle anlatır:
\begin{quote}
  Bir garson masaya tabaklarla tam olarak aynı sayıda bıçak koyduğundan emin
  olmak isterse hiçbirini saymasına gerek yoktur; yeter ki her tabağın sağına
  bir bıçak koysun ve masadaki her bıçak bir tabağın sağında bulunsun. Böylece
  tabaklar ile bıçaklar, üstelik aynı uzamsal ilişki aracılığıyla, birbirleriyle
  biricik biçimde ilişkilendirilmiş olur. \citep[\S70]{Frege1884}
\end{quote}
Bu parçadaki fikir biçimsel bir tanımla açığa çıkarılabilir:

\begin{defn}\ollabel{comparisondef}
  $A$, ancak ve ancak bir !!{bijection} $f\colon A\to B$ varsa $B$ ile
  \emph{eş güçlüdür}; bunu $\cardeq{A}{B}$ ile yazarız.
\end{defn}

\begin{prop}\ollabel{equinumerosityisequi}
Eş güçlülük bir denklik bağıntısıdır.
\end{prop}

\begin{proof}
Eş güçlülüğün yansımalı, simetrik ve geçişli olduğunu göstermeliyiz. $A$,
$B$ ve $C$ kümeler olsun.

\emph{Yansıma.} Bütün $x\in A$ için $\Id{A}(x)=x$ olan özdeşlik eşlemesi
$\Id{A}\colon A\to A$ bir !!{bijection}dur. Dolayısıyla $\cardeq{A}{A}$.

\emph{Simetri.} $\cardeq{A}{B}$ olduğunu, yani bir !!{bijection}
$f\colon A\to B$ bulunduğunu varsayın. $f$ !!{bijective} olduğu için tersi
$f^{-1}$ vardır ve o da !!{bijective}dir. Dolayısıyla $f^{-1}\colon B\to A$
bir !!{bijection}dur; öyleyse $\cardeq{B}{A}$.

\emph{Geçişlilik.} $\cardeq{A}{B}$ ve $\cardeq{B}{C}$ olduğunu, yani
!!{bijection} $f\colon A\to B$ ile $g\colon B\to C$ bulunduğunu varsayın.
O zaman bileşke $\comp{f}{g}\colon A\to C$ !!{bijective}dir; dolayısıyla
$\cardeq{A}{C}$.
\end{proof}

\begin{prop}
$\cardeq{A}{B}$ ise $A$, ancak ve ancak $B$ !!{enumerable} ise
!!{enumerable}dir.
\end{prop}

\begin{editorial}
  Aşağıdaki ispat \olref[enm]{defn:enumerable}'i, \olref[enm]{sec}
  içeriliyorsa; aksi hâlde \olref[enm-alt]{defn:enumerable}'i kullanır.
\end{editorial}

\begin{proof}
$\cardeq{A}{B}$ olduğunu, dolayısıyla bir !!{bijection} $f\colon A\to B$
bulunduğunu; ayrıca $A$'nın !!{enumerable} olduğunu varsayın.
\oliflabeldef{sfr:siz:enm:defn:enumerable}{
  O zaman ya $A=\emptyset$ ya da bir !!{surjective} $g\colon\PosInt\to A$
  fonksiyonu vardır. $A=\emptyset$ ise $B=\emptyset$ de olmalıdır (aksi
  hâlde bir $y\in B$ bulunur, fakat $f(x)=y$ olacak hiçbir $x\in A$
  bulunmazdı). Öte yandan $g\colon\PosInt\to A$ !!{surjective} ise
  $\comp{g}{f}\colon\PosInt\to B$ de !!{surjective}dir. Bunu görmek için
  $y\in B$ olsun. $f$ !!{surjective} olduğundan $f(x)=y$ olacak biçimde bir
  $x\in A$ vardır. $g$ !!{surjective} olduğundan da $g(n)=x$ olacak biçimde
  bir $n\in\PosInt$ vardır. Böylece
\[
(\comp{g}{f})(n)=f(g(n))=f(x)=y
\]
  ve dolayısıyla $\comp{g}{f}$ !!{surjective} olur. $\comp{g}{f}$, $B$'nin
  bir numaralandırmasıdır; öyleyse $B$ !!{enumerable}dir.}
{
  O zaman ya $A=\emptyset$ ya da görüntü kümesi $A$, tanım kümesi ise
  $\Nat$ veya doğal sayıların bir başlangıç parçası olan bir !!{bijection}
  $g$ vardır. $A=\emptyset$ ise $B=\emptyset$ de olmalıdır (aksi hâlde bir
  $y\in B$ bulunur, fakat $f(x)=y$ olacak hiçbir $x\in A$ bulunmazdı).
  Öyleyse söz konusu !!{bijection} $g$'nin bulunduğunu varsayalım. Bu durumda
  $\comp{g}{f}$, görüntü kümesi $B$, tanım kümesi de $g$'ninkiyle aynı (yani
  $\Nat$ ya da onun bir başlangıç parçası) olan bir !!{bijection}dur;
  dolayısıyla $B$ !!{enumerable}dir.}

$B$ !!{enumerable} ise aynı argümanı $f$ yerine !!{bijection}
$f^{-1}\colon B\to A$ ile tekrarlayarak $A$'nın !!{enumerable} olduğunu
elde ederiz.
\end{proof}

\begin{prob}
$\cardeq{A}{C}$ ve $\cardeq{B}{D}$, ayrıca $A\cap B=C\cap D=\emptyset$ ise
$\cardeq{A\cup B}{C\cup D}$ olduğunu gösterin.
\end{prob}

\begin{prob}
$A$ sonsuz ve !!{enumerable} ise $\cardeq{A}{\Nat}$ olduğunu gösterin.
\end{prob}

\end{document}
